\section{Conclusion}

Dans ce travail, nous avons exploré plusieurs modèles de classification pour la détection de plagiat, en nous concentrant particulièrement sur l'évaluation de leur performance à travers des métriques adaptées telles que la précision, le rappel, et la balanced accuracy. Nous avons utilisé des techniques d'agrégation des classificateurs comme le \textit{bagging} et l'\textit{AdaBoost} pour améliorer la robustesse des modèles et leur capacité à généraliser sur des données non vues.

L’analyse des résultats a permis de mettre en lumière l'importance de l'apprentissage dans les performences du modèle. Nous avons aussi pu dégager les modèles plus performants et tester l'impact d'aspects comme le pré-traitement des textes ou la taille des sacs de mots.

De plus, la pré-analyse, en utilisant des outils tels que les nuages de mots, a permis d'identifier les termes récurrents et discriminants, fournissant ainsi un aperçu précieux sur les différences lexicales entre les textes plagiés et non plagiés.

Enfin, des pistes d'amélioration existent l'enrichissement des jeux de données, pour rendre le modèle plus robuste et capable de détecter des cas de plagiat plus subtils.
